% \iffalse meta-comment
%
% Copyright (C) 2017- by Yanshuo Chu <yanshuoc@gmail.com>
%
% This file may be distributed and/or modified under the
% conditions of the LaTeX Project Public License, either version 1.3a
% of this license or (at your option) any later version.
% The latest version of this license is in:
%
% http://www.latex-project.org/lppl.txt
%
% and version 1.3a or later is part of all distributions of LaTeX
% version 2004/10/01 or later.
%
% \fi
%
% \section{学位论文类的第三方依赖}
%
% 第三方宏包的集中加载，一个守卫装完。
%
% \changes{v3.2a}{2026/08/12}{三段 \texttt{*-deps-a/b/c} 合成一段。
%   \texttt{deps-b} 里一行代码都没有，是当年把 \pkg{geometry} 拆出去时留下的空壳}
% \changes{v3.2a}{2026/08/12}{改掉这里写错的顺序理由。原先写「\pkg{geometry} 必须
%   早于本文件里的 \pkg{fancyhdr}」，可 \pkg{fancyhdr} 由排在更前面的
%   \file{hit-thesis-pagestyle.dtx} 加载，本文件那句是重复加载，已删}
% 装配序列上它夹在 \file{hit-thesis-geometry.dtx} 与 \file{hit-thesis-hyperlink.dtx}
% 之间。真正的硬约束只有一条：\pkg{hyperref} 要落到所有包之后，所以
% \file{hit-thesis-hyperlink.dtx} 排在本文件后面。\pkg{geometry} 提前是为了让版心
% 在任何按 \cs{textwidth} 算尺寸的包之前定下来，这条是稳妥起见，不是必须。
% 本文件内部各个包彼此没有先后要求。
%
% ^^A ======================================================================
% ^^A Module thesis-deps: 第三方宏包集中加载与基础配置（含 hyperlink/geometry 内联）（hit-thesis）
% ^^A ======================================================================
%
%
% \subsection{装载宏包}
% \label{sec:loadpackage}
%
% 引用的宏包和相应的定义。
%
% \changes{v3.2a}{2026/08/12}{删掉这里的 \pkg{etoolbox}。
%   \file{src/setup/hit-keylist.dtx} 早就装过一次，这句是重复加载}
%    \begin{macrocode}
%<*thesis-deps>
%    \end{macrocode}
%
% \changes{v3.2a}{2026/08/12}{book-deps-a 的引擎检查改用 \cs{sys\_if\_engine\_xetex:F}，
%   与 art-deps-a 保持一致。报错文本里的空格写成 \texttt{\~{}}}
%    \begin{macrocode}
%    \end{macrocode}
%
% \changes{v3.2a}{2026/08/12}{数学部分抽到 \file{src/utils/hit-math.dtx}}
% 数学宏包、定理机制与数学字体的设置见 \file{src/utils/hit-math.dtx}。
%    \begin{macrocode}
\hit@setup@math
%    \end{macrocode}
%
% 删除 \pkg{newtxtext}，由于最新一次更新 (\texttt{revision 62369}) 中与 \pkg{ctex} 宏包冲突，且无法解决，将正文字体替换为 TG TermesX 与 TG Heros.
% \changes{v3.0.18}{2022/05/02}{删除 \pkg{newtxtext}，由于最新一次更新 (\texttt{revision 62369}) 中与 \pkg{ctex} 宏包冲突，且无法解决，将正文字体替换为 TG Termes 与 TG Heros}
% \changes{v3.0.19}{2022/05/05}{将 TG Termes 字体替换为 TG TermesX 字体，与 \pkg{newtxmath} 宏包更适配，删除 \option{Scale} 选项}
% \changes{v3.2a}{2026/08/12}{英文字体三件套设置从 bookcls 拆出至 book-fonts 模块（与 art-fonts / artplus-fonts 保持一致）}
%    \begin{macrocode}
%    \end{macrocode}
%    \begin{macrocode}
%    \end{macrocode}
%
% 图形支持宏包。
%    \begin{macrocode}
\RequirePackage{graphicx}
%    \end{macrocode}
%
% \changes{v3.2a}{2026/08/12}{原 hithesis.sty 中的 xcolor、tikz 与全局颜色主题并入 book-deps，作为彩色与图形基础宏包}
% 彩色支持与 \pkg{tikz}，封面与正文图形都要。示例插图用的那组颜色跟着它的
% \cs{tikzset} 一起放在 \file{src/hit-tikz-demo.dtx}。
%    \begin{macrocode}
\RequirePackage{xcolor}
\RequirePackage{tikz}
%    \end{macrocode}
%
% \pkg{pdfpages} 宏包便于我们插入扫描后的授权页和声明页 PDF 文档。
%    \begin{macrocode}
\RequirePackage{pdfpages}
\includepdfset{fitpaper=true}
%    \end{macrocode}
%
% 更好的列表环境。
%
% \changes{v3.2a}{2026/08/12}{删掉 \pkg{environ}。它唯一的消费者是四个摘要环境用的
%   \cs{Collect@Body}，那边改用内核的 \texttt{+b} 参数类型之后就没人调了}
%    \begin{macrocode}
\hit@setup@list@packages
%    \end{macrocode}
%
% 禁止 \hologo{LaTeX}自动调整多余的页面底部空白，并保持脚注仍然在底部。
% 脚注按页编号。
%    \begin{macrocode}
\hit@setup@footnote@packages
%    \end{macrocode}
%
% 脚注格式。
%
% \changes{v3.2a}{2026/08/12}{删掉 \option{pifootnote} 与 \pkg{pifont}。带圈脚注号
%   统一由 \cs{quan} 排：先取当前宋体的真字符（\texttt{①}–\texttt{⑩}），超出字体
%   可用范围才画圈。\pkg{pifont} 那套 Dingbats 字形与正文中文字体不搭}
%
% \changes{v3.2a}{2026/08/12}{\pkg{xeCJKfntef} 留在这里。分散对齐的
%   \cs{hit@spread@to} 底层是它的 \env{CJKfilltwosides}，两个类都用，
%   而 \file{src/utils/hit-spread.dtx} 排在 \cs{LoadClass} 之前装不了包}
% \pkg{xeCJKfntef} 的 \env{CJKfilltwosides} 撑开 xeCJK 的 \cs{CJKglue}，
% 封面字段标签的分散对齐靠它，见 \file{src/utils/hit-spread.dtx}。
%
% 表格控制
%    \begin{macrocode}
\RequirePackage{xeCJKfntef}
\RequirePackage{longtable}
%    \end{macrocode}
%
% 使用三线表：\cs{toprule}，\cs{midrule}，\cs{bottomrule}。
%    \begin{macrocode}
\RequirePackage{booktabs}
%    \end{macrocode}
%
% 重定义 \pkg{booktabs} 宏包的默认线宽，绘制三线表时不再需要显式设置，以实现更好的内容与格式分离
% \changes{v3.0.20}{2022/05/06}{重定义 \pkg{booktabs} 宏包的默认线宽，绘制三线表时不再需要显式设置，以实现更好的内容与格式分离}
%    \begin{macrocode}
\hit@setup@table@rules
%    \end{macrocode}
%
% 参考文献引用宏包。
% \changes{v3.0.15}{2020/05/21}{删去natbib package}
%    \begin{macrocode}
% \RequirePackage[sort&compress]{natbib}
%    \end{macrocode}
%
% 生成有书签的 pdf 及其开关，请结合 gbk2uni 避免书签乱码。
% \changes{v3.2a}{2026/08/12}{hyperref/url/gbt7714/citetwocomma 等链接与引用配置从 bookcls 拆出至 book-hyperlink 模块}
%    \begin{macrocode}
%    \end{macrocode}
%    \begin{macrocode}
%    \end{macrocode}
%
% \subsection{页面设置}
% \label{sec:layout}
% 本来这部分应该是最容易设置的，但根据我工\PGR\ 的3.8,3.4,3.2节的版芯矛盾，此处
% 设置两种版芯。
% \changes{v3.2a}{2026/08/12}{页面 geometry 设置从 bookcls 拆出至 book-geometry 模块}
%    \begin{macrocode}
%    \end{macrocode}
%    \begin{macrocode}
%    \end{macrocode}
%
% \changes{v3.2a}{2026/08/12}{book-deps-c 的两处条件加载改用 \cs{legacy\_if:nT}}
% 载入显示行号的包。
% \changes{v1.0.09}{2018/01/07}{添加debug包}
%    \begin{macrocode}
\bool_if:NT \g__hit_debug_bool
  {
    \RequirePackage { layout }
    \RequirePackage { layouts }
    \RequirePackage { lineno }
  }
%    \end{macrocode}
%
% \changes{v3.2a}{2026/08/12}{删掉这里的 \pkg{fancyhdr}。
%   \file{src/flow/hit-thesis-pagestyle.dtx} 在装配序列上排在前面，早就装过了}
% 页眉页脚由 \file{src/flow/hit-thesis-pagestyle.dtx} 加载 \pkg{fancyhdr} 并设置。
%
% 其他包，表格、数学符号包
% \changes{v1.0.06}{2017/10/25}{此处添加子图最后一行图题是否居中选项}
%    \begin{macrocode}
\RequirePackage{tabularx}
\RequirePackage{xltabular}
%    \end{macrocode}
%
% \changes{v3.2a}{2026/08/12}{删掉 \pkg{varwidth}。仓库里它唯一一次出现是建仓那个
%   提交里章标题缩进盒子的注释行，从没真正用过}
%
% \changes{v3.2a}{2026/08/12}{\pkg{changepage} 与 \pkg{multicol} 挪进
%   \file{src/pages/hit-thesis-backmatter.dtx}，只有索引那一页调它们}
% \changes{v3.2a}{2026/08/12}{\pkg{placeins} 与 \pkg{multirow} 挪进
%   \file{src/hit-sty.dtx}。\cs{FloatBarrier} 与 \cs{multirow} 类里一处都不调，
%   是替用户预装的}
%
% \pkg{flafter} 没有对外的宏，装上就重定义内核的 \cs{@addtocurcol}，让浮动体不会
% 出现在引用它的正文之前。扒不出消费者，但类的排版结果依赖它，留在这儿。
%    \begin{macrocode}
\RequirePackage{flafter}
\RequirePackage{subcaption}%支持子图，必须在 ccaption 之前
\RequirePackage{ccaption} %支持双语标题
%    \end{macrocode}
%
% \changes{v3.2a}{2026/08/12}{\cs{bicaption} 支持省略英文图表名的四参数写法}
% \pkg{ccaption} 的 \cs{bicaption} 是五个参数，第四个是第二条题注用的图表名，
% 用户得写成 \texttt{Fig.} 加一个负细空格才能把 \cs{fnum@figure} 插的空格抵掉。
% 这里让第四个参数可以省略，省略时取 \option{const} 族里的 \option{figure-prefix-en}
% 与 \option{table-prefix-en}。
%
% \changes{v3.2a}{2026/08/12}{五参数写法从报错改回可用，只提示一次。
%   v3.1e 的示例正文里就是五参数，报错会让老文档一行都编不过}
% 五个参数的写法照旧可用——那是 \pkg{ccaption} 的标准形式，v3.1e 的示例正文里
% 写的也是它，直接报错会让老文档编不过。现在只提示，v4 再移除。
%    \begin{macrocode}
\hit@setup@caption
\hit@setup@subfigure
%    \end{macrocode}
%
% 中英文索引包。
%    \begin{macrocode}
\RequirePackage[makeindex]{splitidx}
%    \end{macrocode}
%
% \changes{v3.2a}{2026/08/12}{两个索引的声明留在这里，补一条为什么}
% 中英文两个索引。这两句得紧跟 \pkg{splitidx} 的加载：\cs{newindex} 由它定义，
% 而后置那一页（\file{src/pages/hit-thesis-backmatter.dtx}）在装配序列上比
% deps 早，搬过去会撞 Undefined control sequence。试过了。
%    \begin{macrocode}
\newindex[]{china}
\newindex[]{english}
%    \end{macrocode}
%
% 排版logo。
% \changes{v3.0.21}{2022/05/11}{删除 \pkg{xltxtra} 宏包, 来实现规范样式的上标。\issue[125]}
%    \begin{macrocode}
%    \end{macrocode}
%
% \subsection{主文档格式}
% \label{sec:mainbody}
%    \begin{macrocode}
%    \end{macrocode}
%
%
% \subsubsection{字体}
% \label{sec:font}
% \begin{macro}{\normalsize}
% 根据我工规定，正文小四号 (12bp) 字，行距为固定值3--4mm。
%    \begin{macrocode}
%    \end{macrocode}
%
% \end{macro}
% \subsubsection{页眉页脚}
% \label{sec:headerfooter}
% \begin{macro}{\hit@empty}
% \begin{macro}{\hit@plain}
% \begin{macro}{\hit@headings}
% 定义四种页眉页脚格式：
% \begin{itemize}
% \item \texttt{hit@empty}：页眉页脚都没有
% \item \texttt{hit@plain}：只显示页脚的页码。\cs{chapter} 自动调用
% \cs{thispagestyle\{hit@plain\}}。
% \item \texttt{hit@headings}：页眉页脚同时显示
% \item \texttt{hit@scanpage}：只显示页脚的页码，扫描页专用
% \end{itemize}
%
% \end{macro}
% \end{macro}
% \end{macro}
%
% \subsubsection{段落}
% \label{sec:paragraph}
%
% 全文首行缩进 2 字符，标点符号用全角
% \changes{v3.2a}{2026/08/12}{段落相关代码从 bookcls 拆出至 book-paragraph 模块}
%
% \subsubsection{脚注}
% \label{sec:footnote}
% 脚注符合中文习惯，数字带圈。
% \changes{v3.2a}{2026/08/12}{脚注相关代码从 bookcls 拆出至 book-footnote 模块}
%    \begin{macrocode}
%    \end{macrocode}
%
% \subsubsection{数学相关}
% \label{sec:equation}
% 允许太长的公式断行、分页等。
% \changes{v3.2a}{2026/08/12}{数学相关代码从 bookcls 拆出至 book-math 模块}
%    \begin{macrocode}
%</thesis-deps>
%    \end{macrocode}
%
% \subsubsection{浮动对象以及表格}
% \label{sec:float}
% 设置浮动对象和文字之间的距离，由于规范中没有明确规定，根据经验，设置成正文汉字
% 高度。
% \changes{v1.0.09}{2018/01/07}{修正float垂直间距bug}
% \changes{v2.0.09}{2019/06/24}{修正float垂直间距bug}
% \changes{v3.2a}{2026/08/12}{浮动对象与表格代码从 bookcls 拆出至 book-floats 模块；拼接位置留在 cls 中的原处以避免依赖错位}
%
% 按照我工要求，页面中标题之下不少于一行。
%
% \changes{v3.2a}{2026/08/12}{参考文献相关代码从 bookcls 拆出至 book-bib 模块}
%
% \Finale
